%\documentclass[12pt,oneside]{amsart}
\documentclass{scrartcl}

%\documentclass{article}

\usepackage{amsthm,amsmath,amssymb}
\usepackage[numbers]{natbib}

%\usepackage[colorlinks]{hyperref}
\usepackage{hyperref}
\hypersetup{
    colorlinks,%
    citecolor=black,%
    filecolor=black,%
    linkcolor=black,%
    urlcolor=black
}
\usepackage{hypernat}
\usepackage[top=2.5cm, bottom=2.5cm, left=2.8cm,
right=2.8cm]{geometry} 

%\setlength{\parskip}{0pt}
%\setlength{\parsep}{0pt}
%\setlength{\headsep}{0pt}
%\setlength{\topskip}{0pt}
%\setlength{\topmargin}{0pt}
%\setlength{\topsep}{0pt}
%\setlength{\partopsep}{0pt}

\linespread{1}

%\usepackage{marvosym}

%\textwidth = 6.5 in \textheight = 9.2 in \oddsidemargin = 0.0 in
%\evensidemargin = 0.0 in \topmargin = -0.0 in \headheight = 0.0 in
%\headsep = 0.0in \parskip = 0.0in \parindent = 0.0in
%\def\baselinestretch{0.871}


\newtheorem{theorem}{Theorem}[section]
\newtheorem{proposition}[theorem]{Proposition}
\newtheorem{problem}[theorem]{Problem}
\newtheorem{fact}[theorem]{Fact}
\newtheorem{lemma}[theorem]{Lemma}
\newtheorem{defn}[theorem]{Definition}
\newtheorem{corollary}[theorem]{Corollary}
\newtheorem{remark}{Remark}[section]
\newtheorem{example}[theorem]{Example}

\newcommand{\E}{{\mathbb E}}
\newcommand{\se}{{\mathcal{E}}}
\newcommand{\R}{{\mathbb R}}
\renewcommand{\P}{{\mathbb P}}
\newcommand{\ac}{{\mathcal{A}}}
\newcommand{\A}{{\mathcal{A}}}
\newcommand{\B}{{\mathcal{B}}}
\newcommand{\C}{{\mathcal{C}}}
\newcommand{\M}{{\mathcal{M}}}
\newcommand{\Mf}{{\mathfrak{M}}}
\newcommand{\T}{{\mathcal{T}}}
\newcommand{\Z}{{\mathcal{Z}}}
\newcommand{\K}{{\mathcal{K}}}
\newcommand{\lc}{{\mathcal{L}}}
\newcommand{\Ps}{{\mathcal{P}}}
\newcommand{\G}{{\mathcal{G}}}
\newcommand{\pa}{{\mathring{p}}}
\newcommand{\F}{{\cal F}}
\newcommand{\samp}{{\mathcal{X}}}
\newcommand{\X}{{\mathcal{X}}}
\newcommand{\hi}{{\hat{\Phi}}}
\newcommand{\data}{{\text{data}}}

\newcounter{rcnt}[section]
\renewcommand{\thercnt}{(\roman{rcnt})}

\def\qt#1{\qquad\text{#1}}

\def\argmin{\mathop{\rm argmin}}
\def\argmax{\mathop{\rm argmax}}


\setlength{\parskip}{1.4 \medskipamount}
%\sloppy
%\linespread{1.3}

\begin{document}

\title{STAT 238 - Bayesian Statistics  \\
  Lab Two} 
% \subtitle{Spring 2026, UC Berkeley}

\author{Aditya Guntuboyina}


\date{30 Jan 2026}

\maketitle

\section{Uniform Prior on $\log \sigma$ or $\sigma$}

Consider the problem of estimating a scale parameter $\sigma$ from
observations $X_1, \dots, X_n \overset{\text{i.i.d}}{\sim} N(0,
\sigma^2)$. We want to use an uniformative prior for $\sigma$. There
are two options. Option A is $\sigma \sim \text{uniform}(0, +\infty)$
(or $\text{uniform}(0, C)$ for a large $C$). Option B is $\log \sigma
\sim \text{uniform}(-\infty, +\infty)$ (or $\text{uniform}(-C, C)$ for
a large $C$). The second prior (option B) is universally preferred to
the first prior (option A). Here are some reasons for why this is the
case.

\begin{enumerate}
\item \textbf{Relative Scale vs Absolute Scale}: Consider the two
  probabilities $\P\{1 < \sigma < 2\}$ and $\P\{101 < \sigma <
  102\}$. Under the first prior (uniform$(0, C)$), both these
  probabilities are the same.

  Under the second prior ($\log \sigma$ is
  uniform on $(-C, C)$), we have:
  \begin{align*}
    \P\{1 < \sigma < 2\} = \P\{0 < \log \sigma < \log 2\} = \frac{\log
    2}{2C} = \frac{0.693}{2C}
  \end{align*}
  and
  \begin{align*}
    \P\{101 < \sigma < 102\} = \P\{\log 101 < \log \sigma < \log 102\}
    = \frac{\log(102/101)}{2C} = \frac{0.00985}{2C} 
  \end{align*}
  The second probability is therefore much smaller than the
  first. This reflects the fact that Prior B regards the values 101
  and 102 as much closer to one another than the values 1 and 2. This
  behavior aligns well with intuition on a relative scale: moving from
  $\sigma = 1$ to $\sigma = 2$ corresponds to a doubling, whereas
  moving from $\sigma = 101$ to $\sigma = 102$ represents only a very
  small relative change.

\item \textbf{Prior Odds}: Under the first prior, consider the prior
  odds that $\sigma < 2$ 
  versus $\sigma \geq 2$:
  \begin{align*}
    \frac{\P\{\sigma < 2\}}{\P\{\sigma \geq 2\}} = \frac{2/C}{(C-2)/C}
    = \frac{2}{C - 2}
  \end{align*}
  which is very small (as $C$ is large). This is clearly an
  informative statement about $\sigma$ (that it is much more likely to
  be more than 2 than smaller than 2). On the other hand, the same
  odds for the second prior becomes:
  \begin{align*}
    \frac{\P\{\sigma < 2\}}{\P\{\sigma \geq 2\}} = \frac{\P\{\log \sigma <
    \log 2\}}{\P\{\log \sigma \geq \log 2\}} = \frac{C + \log 2}{C -
    \log 2}
  \end{align*}
  which is approximately equal to 1, reflecting prior ignorance
  between the events $\sigma < 2$ and $\sigma \geq 2$.

\item \textbf{Posterior Propriety for $n = 1$}: The posterior for
  $\sigma$ corresponding to the first prior is: 
  \begin{align}\label{post1}
    \frac{2 \left(\frac{S}{2}
    \right)^{(n-1)/2}}{\Gamma\left(\frac{n-1}{2} \right)}
    \sigma^{-n} \exp \left(-\frac{S}{2 \sigma^2} \right) I\{\sigma > 0\}
  \end{align}
  while the posterior corresponding to the second prior is:
  \begin{align}
    \label{post2}
    \frac{2 \left(\frac{S}{2}
    \right)^{n/2}}{\Gamma\left(\frac{n}{2} \right)}
    \sigma^{-(n+1)} \exp \left(-\frac{S}{2 \sigma^2} \right) I\{\sigma >
    0\}. 
  \end{align}
  In both these formulae, $S := \sum_{i=1}^n X_i^2$. 
  
  These two posteriors will behave similarly when $n$ is
  large. However, when $n = 1$, the first posterior is actually
  ill-defined because the Gamma function term is $\Gamma(0)$ which is
  $\infty$. In other words, the posterior is improper when $n = 1$.
  On the other hand, the second posterior is proper even when $n =
  1$.

  This is another reason for preferring the second prior in this
  problem. We would intuitively expect to obtain some concrete
  information on $\sigma$ when $n = 1$ so we want the posterior to be
  proper in that case. This is only true for the second prior but not
  for the first prior. 
\end{enumerate}

Now consider the problem of estimating both $\theta$ and $\sigma$ from
$X_1, \dots, X_n \overset{\text{i.i.d}}{\sim} N(\theta,
\sigma^2)$. Here there are two options for the prior. Prior A is:
\begin{align*}
  \theta, \sigma \text{ are independent with } \theta \sim
  \text{unif}(-\infty, \infty) \text{ and } \sigma \sim \text{unif}(0, \infty). 
\end{align*}
Prior B is
\begin{align*}
  \theta, \log \sigma \overset{\text{i.i.d}}{\sim} \text{unif}(-\infty,
  \infty). 
\end{align*}
Again the universally preferred prior is the second one. All of the
reasons previously mentioned apply here as well. For the third point, 
the marginal posterior of $\sigma$ corresponding to the first prior
is:
\begin{align}
  \label{post11}
      \frac{2 \left(\frac{S}{2}
    \right)^{(n-2)/2}}{\Gamma\left(\frac{n-2}{2} \right)}
    \sigma^{-(n-1)} \exp \left(-\frac{S}{2 \sigma^2} \right) I\{\sigma > 0\}
\end{align}
and the marginal posterior for $\sigma$ corresponding to the second
prior is:
\begin{align}
  \label{post22}
  \frac{2 \left(\frac{S}{2}
    \right)^{(n-1)/2}}{\Gamma\left(\frac{n-1}{2} \right)}
    \sigma^{-n} \exp \left(-\frac{S}{2 \sigma^2} \right) I\{\sigma >
  0\}. 
\end{align}
In both the above formulae, $S = \sum_{i=1}^n (X_i - \bar{X})^2$. The
first posterior above is ill-defined when $n = 1, 2$  while the second
posterior is ill-defined only for $n = 1$. It is well-known that in
this problem we would need at least two observations to estimate
$\sigma$, so we would like the posterior to be well-defined for $n =
2$ which is only true if we use the second prior. 

For more, see \citet[pages 41-47]{zellner}. See also \citet[Section 12.4]{jaynes}
for another justification for the $\log \sigma \sim
\text{uniform}(-\infty, +\infty)$ prior based on invariance to certain
parameter transformations. 


\begin{thebibliography}{9}

\bibitem[Jaynes(2003)]{jaynes} Jayes, E. T. (2003)
\textit{Probability theory: the logic of science}. Cambridge
University Press, 2003.
  
\bibitem[Zellner(1971)]{zellner} Zellner, A., (1971)
\textit{An introduction to Bayesian inference in
Econometrics}. Wiley Classics Library, 1971 (reprint edition in
1996). 
  
  \end{thebibliography}



\end{document}

%%% Local Variables:
%%% mode: latex
%%% TeX-master: t
%%% End:
